\section{Integral normal forms for the constants on a marked arc}

% BEGIN A82N1
\paragraph{A82N1. Reading the two numerator conditions in powers of the level factors.}
Use an arc of I82B2--I82B3 in backward order $s=0,\ldots,\ell-1$.
Put $A=O[[v]]$, let $L$ be the coefficient field and $\lambda$ a
uniformizer. In L81S2's tame invariant coordinate let
$P_s=\varphi^s(v+p)$. It has distinguished degree $p^s$:
up to a unit $v$ is $X^{p-1}$, the substitution by Frobenius
modulo $p$ raises the original variable to the $p$th power, and
$P_s(0)=p$. Its roots correspond to the tame orbits of the
primitive $p^{s+1}$ roots in $1+X$. Indeed the linearizing
coordinate $Z$ of A81R3 has an integral series inverse, and
$Z^{p-1}$ separates its nonzero tame orbits. Apply this to
$\varphi^s(v+p)$ using the level-one root of $v+p$ from L81S2.
Consequently the different $P_s$ are relatively prime over $L$,
and Weierstrass preparation gives unit multiples of monic
distinguished polynomials in $A$. Quotienting by any power gives
a free $O$-module on the powers of $v$ below the polynomial
degree. Coefficient denominators can therefore be cancelled in
divisibility by any product of these factors.

Use constant coordinates $T=O^r$ at the initial cut and let
$U_s$ be its intersection with a specified rational subspace
in $T_L=T[1/p]$. In particular these may be the transported
$\mathrm{Fil}^{b_s}$ at an up position and $\mathrm{Fil}^{-\eta}$
at a down position in the original rounded comparison.
They are saturated direct summands, since their quotients are
torsion free over the DVR. For the given exponents $H_s$ let
$\Lambda=\prod_sP_s^{H_s}$, with $H_s\ge b_s$ at an up
position. The integral numerator $y\in A\otimes T$ represents
$x=y/\Lambda$. At level $s$ take the Taylor expansions at
\emph{all} roots of $P_s$ after any splitting extension. The
rounded homogeneous conditions on $x$ are respectively
\[
 \begin{array}{c|l}
 \text{thin}&x\text{ regular}\\
 \text{up}&P_s^{b_s}x\text{ regular and }x\bmod U_s
                                        \text{ regular}\\
 \text{down}&x\text{ regular and }x\bmod U_s
                                      \equiv0\pmod{P_s^{b_s}} .
 \end{array}
\]
Reduction modulo $U_s$ means projection to $(T/U_s)_L$.
Write $\mathcal N$ for integral numerators satisfying this table.
Write $\mathcal I$ for the same conditions allowing subtraction
of one vector of $T_L$ independently at each level, common to
the roots of that level. On the thin and up rows adding such a
constant preserves both directions of the test. On the down
row the projected residue may instead be a constant.

These are exactly the source's algebraic $\mathrm{NN}$ and
$\mathrm{II}$. Indeed after inverting $\lambda$, the other factors
are units modulo each power of $P_s$, so divisibility can be
tested on these jets, by polynomial remaindering or Chinese
remaindering. In coordinates adapted to $U_s$ the conditions
on $y$ are membership in $P_s^{H_s}A_L^r$ on a thin row, in
$P_s^{H_s-b_s}(U_sA_L)+P_s^{H_s}A_L^r$ on an up row,
and in $P_s^{H_s}(U_sA_L)+P_s^{H_s+b_s}A_L^r$ on a down
row, where $A_L=A[1/p]$. Clearing finitely many coefficient
denominators is precisely the uniformizer power in
$\mathrm{LevelCond}$. The same clearing applied to the subtracted
constant gives its integral representative in $\mathrm{II}$.
This identifies the normal forms below with the numerator tests
in the framed level comparison. Constructing the frame and lifting
its horizontal sections are additional parts of the physical arc
problem.
% END A82N1

% BEGIN A82N2
\paragraph{A82N2. No down position and one down position.}
If there is no down position, A82N1 directly gives
$\mathcal I=\mathcal N$, allowing either one or two up positions.
With a single down position and all other positions thin, every
$y\in\mathcal I$ has the form $y=\Lambda d+\nu$ with $d\in T$
and $\nu\in\mathcal N$. In fact all jets of $x=y/\Lambda$ are
regular, hence $x$ is integral over the power series ring in the
explicit sense $x\in A\otimes T$: each $P_s^{H_s}$ divides $y$
rationally and the product divisibility is integral by A82N1.
Let $D=P_d^{b_d}$ at the down position $d$. The residue of $x$
in $(T/U_d)\otimes A/(D)$ is a rational constant, which is
integral since $1$ is part of the integral polynomial basis of
$A/(D)$. Lift it to an integral constant in $T$ and subtract.
This gives the form claimed. Here and below every thick width
$b_d$ is positive by I82W3. The reverse inclusion follows by
permitting constants at the down row.

The same one-constant answer holds when one up position $u$
precedes the down position $d$ in the backward list. Put
$P=P_u^{b_u}$ and $D=P_d^{b_d}$. Now $q=Px\in A\otimes T$
by the same divisibility argument, and
$q\in U_u A+P(A\otimes T)$. Modulo $D$, its projection onto
$T/U_d$ is $P$ times a rational constant $z$. The size condition
of I82B3 is $b_u<pb_d$, hence $p^u b_u<p^d b_d$. In the
integral free quotient $A/(D)$ the reduction of $P$ modulo
$\lambda$ is $v^{p^u b_u}$ times a unit, nonzero by this
strict inequality. Some coefficient of its remainder is a unit.
The integrality of $Pz$ therefore forces $z$ integral, by
applying an $O$-linear functional on the quotient taking $P$
to one. Lift $z$ to $T$ and subtract it from $x$ as before.
At the up position such subtraction preserves the homogeneous
condition. This gives the original no-down, single-down and
up-before-down normal forms with integral constants.
% END A82N2

% BEGIN A82N3
\paragraph{A82N3. Two down positions and integral interpolation.}
For two down positions write $D_j=P_j^{b_j}$ at these positions,
indexed here by $j=0,1$ in their arc order. Let $J=(D_1,D_0)$.
Its intersection with $O$ is $(\lambda^\rho)$ for some finite
$\rho$, since polynomial Bezout after inverting $\lambda$ gives
a nonzero constant in $J$. Choose $f\in(D_1)$ with
$f-\lambda^\rho\in(D_0)$ using this equality.
An integral numerator in $\mathcal I$ has precisely the form
\[
              y=\Lambda d+\Lambda f d'+\nu,
              \qquad d,d'\in T,\quad\nu\in\mathcal N .
\]
To see this, all jets of $x=y/\Lambda$ are regular, so $x$
is integral. At the two down positions its projected residues
are integral constants as in A82N2. Lift them to $d_0,d_1\in T$.
Thus $d_0-d_1\in(U_0+U_1)A+J(A\otimes T)$. This implies
$d_0-d_1\in U_0+U_1+\lambda^\rho T$ even if the sum of the
two saturated submodules is itself unsaturated. Indeed take
Smith coordinates for that sum. In a coordinate with Smith
ideal $(\lambda^k)$, the constants in $(\lambda^k)+J$ form
$(\lambda^{\min(k,\rho)})$: if $\lambda^i=\lambda^k z+j$
with $i<k$ and $j\in J$, the factor $1-\lambda^{k-i}z$ is a
unit, so $\lambda^i\in J$. A zero Smith ideal gives $J\cap O$
directly. The claim follows coordinate by coordinate.
Choose $d\equiv d_1\bmod U_1$ and $d'$ such that
$d+\lambda^\rho d'\equiv d_0\bmod U_0$ by that inclusion.
Subtract $d+fd'$ from $x$. At each down position the homogeneous
test now holds, and all other positions are thin. Conversely the
two displayed terms have projected residues that are constants
at both down positions. This proves both inclusions with the
minimal integral interpolation power $\lambda^\rho$.
% END A82N3

% BEGIN A82N4
\paragraph{A82N4. A down position before an up position.}
Let $D=P_d^{b_d}$ and $P=P_u^{b_u}$ with $d<u$. Write
$P=\lambda^c\alpha+D\beta$, where $c$ is the minimum valuation
of the coefficients of the remainder of $P$ in $A/(D)$.
It is finite since the levels are relatively prime. Thus the
image of $\alpha$ in $A/(D)$ is primitive and admits an
$O$-linear functional taking it to one. Choose the integral
$\alpha,\beta$ by lifting the divided remainder. Then the
complete integral form is
\[
  y=\Lambda d_0+\frac{\Lambda}{P}\alpha d_1+\nu,
  \qquad d_0\in T,\quad d_1\in U_u,\quad\nu\in\mathcal N .
\]
Indeed $q=Px$ is integral as in A82N2, with
$q\in U_u A+P(A\otimes T)$. At the down position
$q\equiv Pz\bmod U_d$ in $(A/(D))[1/p]$, for a rational
constant $z\in(T/U_d)_L$. Applying the functional on $\alpha$
shows that $\lambda^c z$ belongs to the integral quotient $T/U_d$.
Lift it to $z_0\in T$. The two congruences for $q$ now give,
in $(A/(D))\otimes T$,
\[
       \alpha z_0\in (U_u+U_d)(A/(D))
                           +\lambda^c ((A/(D))\otimes T).
\]
Applying the same functional shows
$z_0\in U_u+U_d+\lambda^c T$. Choose $d_1\in U_u$ and
$d_0\in T$ with $z_0\equiv d_1+\lambda^c d_0\bmod U_d$.
At the down level $\alpha/P=\lambda^{-c}$ modulo $D$ after
inverting coefficients. Thus subtracting
$d_0+(\alpha/P)d_1$ from $x$ kills its projected residue there.
At the up level the residue of the pole is in $U_u$, as required,
and the remaining levels are regular. This proves membership
of the remainder in $\mathcal N$. Conversely the displayed
numerator is integral since $P\mid\Lambda$ and the same local
checks give $\mathcal I$. This proof keeps the exact coefficient
power even if $U_u+U_d$ is unsaturated.
% END A82N4

% BEGIN A82N5
\paragraph{A82N5. The lifting cost of the constants.}
Use the two numerator models for the source and target of the
tested lattice map. Their map $f_0:T_1\to T_2$ on constant
coordinates at the preceding cut is evaluation of the Wach map
at zero. R82W4 gives $p^{rr}T_2\subset f_0T_1$ there: lift
a vector to the target Wach module at that cut, lift its scaled
value through the Wach map, and reduce modulo the invariant
coordinate $v$. Use tame averaging for invariant representatives.
The basis of the invariant lattice specializes to an integral
basis by the argument of L81S2. This proves the original
$\mathrm{LLift}$ with $rr=rr_{\rm par}$. The constant matrix
extends coefficientwise to $A$.

For any target $y\in\mathcal I_2$, write its normal form in
A82N2--A82N4 with remainder $\nu\in\mathcal N_2$.
On a no-down arc one may take source numerator zero. For a
single down, up-before-down or double down, lift each displayed
constant after multiplication by $p^{2a_0}$, since the native
$a_0$ has $2a_0\ge rr$. Use the same scalar functions of the
level polynomials as in the target normal form. The resulting
integral source numerator $y_1$ lies in $\mathcal I_1$ and
\[
   p^{2a_0}y=(1\otimes f_0)y_1+p^{2a_0}\nu .
\]
In the down-before-up case the ordinary constant lifts in the same
way; the second one must lift inside the selected up sublattice.
Precisely the additional sufficient input is
$p^{2a_0}U_{2,u}\subset f_0(U_{1,u})$. Under that input the
display follows from A82N4 as well. These are the integral
$\mathrm{ConstLift}$ conclusions for the normalized numerator
models. Identification with maps on horizontal sections requires
the corresponding localization and frame comparisons.
% END A82N5
